%% ============================================================
%% Chapter 1 — Introduction
%% ============================================================
\chapter{Introduction}
\label{chap:introduction}

\section{Motivation}
\label{sec:intro_motivation}

The proliferation of large-scale neural language models has enabled Question
Answering (QA) systems of remarkable fluency. Modern systems based on
Transformer architectures \cite{devlin2019bert,liu2019roberta} can produce
grammatically well-formed, contextually plausible responses even in the absence
of genuine evidence. This phenomenon—commonly called \emph{hallucination}—is
particularly acute in \emph{Retrieval-Augmented Generation} (RAG) pipelines
\cite{lewis2020rag}, where a retrieval module first fetches candidate passages
from a corpus and a neural reader then extracts or synthesises an answer.

When the retrieval step returns irrelevant or empty passages, the reader
operates on a flawed context. Rather than signalling uncertainty, it often
fabricates an answer—a failure mode with potentially severe consequences in
high-stakes domains such as medicine, law, and education. Consider a clinical
decision-support system queried about a rare drug interaction: an overconfident
but wrong answer can mislead a clinician, while an honest ``I do not know''
prompts the user to consult authoritative sources.

The core tension is therefore between \emph{coverage} (answering as many
questions as possible) and \emph{precision} (ensuring that answers provided are
correct). An ideal system navigates this trade-off gracefully by choosing to
\emph{abstain}—withhold an answer—precisely when its internal evidence is
insufficient. This capacity for principled abstention is the focus of the
present work.

\section{Research Question}
\label{sec:intro_rq}

We ask:

\begin{quote}
  \textit{How can a Retrieval-Augmented QA pipeline accurately detect
  unanswerable questions and abstain, thereby reducing false answers while
  preserving useful coverage on answerable questions?}
\end{quote}

Concretely, we investigate whether signals derived from BM25 retrieval scores
and reader confidence scores can—individually or jointly—serve as reliable
proxies for answerability, and how a lightweight classifier combining these
features compares to simpler threshold-based heuristics.

\section{Contributions}
\label{sec:intro_contributions}

This work makes the following contributions:

\begin{enumerate}
  \item \textbf{Systematic feature engineering for answerability.}  We identify
        and extract ten inference-available features from BM25 retrieval and RoBERTa
        reader outputs that jointly characterise the likelihood that a question
        is answerable from the retrieved context.

  \item \textbf{Three-strategy comparison.}  We design and evaluate three
        abstention strategies of increasing complexity—a \emph{non-abstaining
        baseline}, a \emph{BM25 score threshold}, and a \emph{logistic
        regression answerability classifier}—on a held-out test partition of
        SQuAD~2.0.

  \item \textbf{Rigorous experimental protocol.}  We introduce an
        \emph{article-level} train/validation/test split to prevent context
        leakage, tune all hyper-parameters exclusively on the validation set,
        and evaluate on a fully unseen test set.

  \item \textbf{Error taxonomy.}  We categorise system errors into retrieval
        failures, extraction failures, and classification failures, providing
        qualitative examples that illuminate the remaining challenges.
\end{enumerate}

\section{Outline of the Three Abstention Strategies}
\label{sec:intro_strategies}

\paragraph{Strategy 1 — Always Answer (Baseline).}
The system always passes the top-retrieved passage to the reader and returns
whatever answer it produces, without any abstention mechanism. This represents
standard RAG behaviour and sets the precision/coverage upper and lower bounds,
respectively.

\paragraph{Strategy 2 — BM25 Threshold.}
If the BM25 score of the top-retrieved passage exceeds a threshold $\tau_{\text{BM25}}$
(tuned on the validation set by maximising F1), the system answers; otherwise
it abstains. This strategy leverages a single, interpretable retrieval signal.

\paragraph{Strategy 3 — Answerability Classifier.}
A logistic regression model is trained on ten inference-available features drawn from BM25
retrieval and reader outputs. It predicts a probability of answerability; if
this probability exceeds a threshold (again tuned on validation), the system
answers. This strategy combines complementary signals and can capture
non-linear interactions among them.

\section{Report Structure}
\label{sec:intro_structure}

The remainder of this report is organised as follows.
Chapter~\ref{chap:background} reviews the relevant literature on reading
comprehension datasets, sparse retrieval, Transformer-based QA, and selective
prediction.
Chapter~\ref{chap:methodology} details the experimental methodology, including
dataset preparation, feature extraction, and strategy definitions.
Chapter~\ref{chap:architecture} presents the end-to-end system architecture.
Chapter~\ref{chap:experiments} describes the experimental setup and evaluation
protocol.
Chapter~\ref{chap:results} reports quantitative results and discusses figure
evidence.
Chapter~\ref{chap:error_analysis} provides a qualitative error analysis.
Chapter~\ref{chap:discussion} situates findings in the broader research context.
Chapter~\ref{chap:conclusion} concludes and outlines directions for future work.
